% Options for packages loaded elsewhere
\PassOptionsToPackage{unicode}{hyperref}
\PassOptionsToPackage{hyphens}{url}
\documentclass[
  12pt,
]{article}
\usepackage{xcolor}
\usepackage[margin=1in]{geometry}
\usepackage{amsmath,amssymb}
\setcounter{secnumdepth}{-\maxdimen} % remove section numbering
\usepackage{iftex}
\ifPDFTeX
  \usepackage[T1]{fontenc}
  \usepackage[utf8]{inputenc}
  \usepackage{textcomp} % provide euro and other symbols
\else % if luatex or xetex
  \usepackage{unicode-math} % this also loads fontspec
  \defaultfontfeatures{Scale=MatchLowercase}
  \defaultfontfeatures[\rmfamily]{Ligatures=TeX,Scale=1}
\fi
\usepackage{lmodern}
\ifPDFTeX\else
  % xetex/luatex font selection
  \setmainfont[]{Times New Roman}
\fi
% Use upquote if available, for straight quotes in verbatim environments
\IfFileExists{upquote.sty}{\usepackage{upquote}}{}
\IfFileExists{microtype.sty}{% use microtype if available
  \usepackage[]{microtype}
  \UseMicrotypeSet[protrusion]{basicmath} % disable protrusion for tt fonts
}{}
\makeatletter
\@ifundefined{KOMAClassName}{% if non-KOMA class
  \IfFileExists{parskip.sty}{%
    \usepackage{parskip}
  }{% else
    \setlength{\parindent}{0pt}
    \setlength{\parskip}{6pt plus 2pt minus 1pt}}
}{% if KOMA class
  \KOMAoptions{parskip=half}}
\makeatother
\usepackage{graphicx}
\makeatletter
\newsavebox\pandoc@box
\newcommand*\pandocbounded[1]{% scales image to fit in text height/width
  \sbox\pandoc@box{#1}%
  \Gscale@div\@tempa{\textheight}{\dimexpr\ht\pandoc@box+\dp\pandoc@box\relax}%
  \Gscale@div\@tempb{\linewidth}{\wd\pandoc@box}%
  \ifdim\@tempb\p@<\@tempa\p@\let\@tempa\@tempb\fi% select the smaller of both
  \ifdim\@tempa\p@<\p@\scalebox{\@tempa}{\usebox\pandoc@box}%
  \else\usebox{\pandoc@box}%
  \fi%
}
% Set default figure placement to htbp
\def\fps@figure{htbp}
\makeatother
\setlength{\emergencystretch}{3em} % prevent overfull lines
\providecommand{\tightlist}{%
  \setlength{\itemsep}{0pt}\setlength{\parskip}{0pt}}
\usepackage{bookmark}
\IfFileExists{xurl.sty}{\usepackage{xurl}}{} % add URL line breaks if available
\urlstyle{same}
\hypersetup{
  pdftitle={Implementazione Parallela dell'Automa Cellulare Game of Life in C++ e OpenMP},
  pdfauthor={Elias Ezequiel Moreno matricola n° 244211},
  hidelinks,
  pdfcreator={LaTeX via pandoc}}

\title{Implementazione Parallela dell'Automa Cellulare Game of Life in
C++ e OpenMP}
\author{Elias Ezequiel Moreno matricola n° 244211}
\date{}

\begin{document}
\maketitle

Il gioco della Vita di John Conway è un modello di automa celllulare.
Questo ``gioco'' si svolge in una griglia quadrata NxN, composta da
celle che possono assumere solo due stati: 0(la cellula è morta) o 1(la
cellula è viva). Lo stato della griglia evolve in un tempo discreto, con
una separazione netta tra lo stato precedente e quello futuro. Il valore
di ogni cella varia a seconda dello stato attuale delle celle adiacenti
ad essa, seguendo le seguenti regole:

\begin{itemize}
\tightlist
\item
  Qualsiasi cella viva con meno di due celle vive adiacenti muore, come
  per effetto d'isolamento;
\item
  Qualsiasi cella viva con due o tre celle vive adiacenti sopravvive
  alla generazione successiva;
\item
  Qualsiasi cella viva con più di tre celle vive adiacenti muore, come
  per effetto di sovrappopolazione;
\item
  Qualsiasi cella morta con esattamente tre celle vive adiacenti diventa
  una cella viva, come per effetto di riproduzione.
\end{itemize}

\subsubsection{Studio delle performance al variare della dimensione del
problema da
risolvere}\label{studio-delle-performance-al-variare-della-dimensione-del-problema-da-risolvere}

L'analisi delle perfomance viene svolta paragonando i risultati ricavati
dal codice sequenziale con quello in parallelo, mantenendo il numero di
thread pari al numero di core presenti nel dispositivo e le iterazioni
pari a 100:

\pandocbounded{\includegraphics[keepaspectratio]{Implementazione-Parallela_files/figure-latex/unnamed-chunk-3-1.pdf}}

Da quello che si può osservare dal grafico, il parallelismo risulta
essere uno strumento molto potente, in grado di velocizzare il codice
soprattutto aumentando il livello di complessità della matrice come si
può notare

\subsubsection{Studio della variazione delle performance in base alla
variazione del
parallelismo}\label{studio-della-variazione-delle-performance-in-base-alla-variazione-del-parallelismo}

\paragraph{Strong Scaling}\label{strong-scaling}

Le perfomance del nostro codice possono variare non solo a come le
istruzioni vengono processate (sequenziale o parallelo) ma anche in base
a come esse siano distribuite. Questo punto è in grado di essere
regolato in base al numero di thread che noi definiamo da usare.
Aumentare il numero di thread non comporta sempre l'aumento delle
perfomance dato che, se aumentati oltre un certo limite può causare un
fenomeno chiamato overhead di gestione, dove l'elevato numero di thread
genera una contesa nelle risorse del computer, provocando una riduzione
nell'efficienza.

Svolgo il test con una matrice 10000x10000, tenendo conto che in
sequenziale il tempo di esecuzione è di 4.99 sec

\pandocbounded{\includegraphics[keepaspectratio]{Implementazione-Parallela_files/figure-latex/unnamed-chunk-5-1.pdf}}

Il grafico presenta un comportamento estremamente peculiare e
identificativo delle caratterstiche del nostro hardware. Infatti per un
numero di thread pari al numero di core del nostro computer si ottiene
lo speedup maggiore in caso in cui si utilizzi come ottimizzatore
``-O3'', per un numero di thread pari a 16, sono gli altri due
ottimizzatori che ottengono il maggior speedup, coincidendo con il
numero di processori che ha il computer. Aumentare il numero di thread
ulteriormente risulta controproducente, questo dovuto ai ritardi
generati dalla contesa delle risorse hardware e dall'alto overhead.

\pandocbounded{\includegraphics[keepaspectratio]{Implementazione-Parallela_files/figure-latex/unnamed-chunk-6-1.pdf}}

Nel grafico dell'efficienza si può notare invece, come tale si riduce
drasticamente man mano che il numero di thread aumenta fino a quando
tale valore non risulta quasi zero, dimostrando come aumentare il numero
di thread all'infinito non porti un beneficio nelle performance del
codice.

\paragraph{Weak Scaling}\label{weak-scaling}

facendo un analisi dello weak scaling, vogliamo verificare l'efficienza
di ciascun thread, aumentando il numero di thread ma mantenendo fisso il
carico di lavoro per ciascuno di esso. Questo metodo di analisi va a
confrontare il tempo di esecuzione che, se il programma è debolmente
scalabile, rimarrà costante circa, altrimenti esso andrà poco a poco
aumentando. Si può dedurre che il tempo aumenterà sempre più a causa
della \textbf{Banda passante della Memoria} ed all'\textbf{Overhead di
Sincronizzazione}.

\pandocbounded{\includegraphics[keepaspectratio]{Implementazione-Parallela_files/figure-latex/unnamed-chunk-8-1.pdf}}

\subsubsection{Studio della variazione delle performance in base al
numero di variabili condivise/private, e in base al modo di scheduling
dei
threads}\label{studio-della-variazione-delle-performance-in-base-al-numero-di-variabili-condiviseprivate-e-in-base-al-modo-di-scheduling-dei-threads}

\paragraph{Private vs Shared}\label{private-vs-shared}

In un ambiente di calcolo parallelo a memoria condivisa, la corretta
classificazione delle variabili è fondamentale per garantire sia la
correttezza matematica della simulazione, sia la sua efficienza. In
OpenMp, la memoria viene gestita attraverso due classi principali:

\begin{itemize}
\tightlist
\item
  \textbf{Condivisa}(shared): sono le variabili a cui posso avere
  accesso tutti i thread attivi contemporaneamente. Nel nostro codice,
  le variabili condivise sono la dimensione N, e gli stati CurrentGen e
  NextGen, dato che ogni thread elabora celle strettamente diverse, che
  quindi non vanno in conflitto. Usare variabili condivise permette di
  evitare che inutili copie vengano fatte durante l'esecuzione del
  programma, evitando quindi il più possibile i tempi di scrittura e
  lettura, ossia i più lenti.
\item
  \textbf{Privata}(Private): sono le variabili che ciascun thread ha, il
  cui valore dev'essere indipendente dall'azione degli altri thread.
  Affinché l'algoritmo funzioni correttamente, indici di scorrimento
  come x,ydx e dy, e le variabili temporanee come Alive e CurrentState,
  devono essere private. Questo perché se fossero condivise, ogni thread
  sarebbe in grado di sovrascrivere tale valore, portando a perdite di
  informazioni a ``salti'' o sovrascritture della variabile, causando
  quindi un'avanzamento errato dell'algoritmo che porta ad un risultato
  non concorde
\end{itemize}

Una prova molto interessante da compiere in modo da capire cosa potrebbe
accadere in caso di un'errata configurazioni di variabili condividise e
private è quella di forzare il codice a identificare CurrentGen come
variabile privata. Quello che potremmo dedurre che accada è che il
programma andrà a generare una griglia per ciascun thread, aumentando
non il tempo di lettura.

\pandocbounded{\includegraphics[keepaspectratio]{Implementazione-Parallela_files/figure-latex/unnamed-chunk-10-1.pdf}}

\paragraph{Static vs Dynamic vs
Guided}\label{static-vs-dynamic-vs-guided}

Lo scheduling è un altro requisito importante per l'efficienza di un
programma. Esso è il processo mediante il quale un sistema decide
l'ordine di esecuzione, assegnando le risorse necessarie. Si divide
principalmente in due approcci, che si differenziano in base al momento
in cui vengono fatte tali decisioni:

\begin{itemize}
\tightlist
\item
  \textbf{Statico}(Static Scheduling): l'assegnazione delle risorse
  viene decisa prima dell'esecuzione del programma. Questo permette di
  migliorare l'overhead ma manca assolutamente di flessibilità.
\item
  \textbf{Dinamico}(Dynamic Scheduling): le risorse vengono assegnazte
  durante l'esecuzione del programma, permettendo quindi maggiore
  flessibilità in caso di variazione del carico di lavoro, aumentando il
  tempo di esecuzione del programma.
\item
  \textbf{Guidato}(Guided Scheduling): le risorse vengono assegnate
  durante l'esecuzione in blocchi (chunk) di dimensioni via via
  decrescenti fino a un limite minimo prefissato, permettendo di
  mantenere l'alta flessibilità ai carichi di lavoro variabili ma
  riducendo drasticamente il tempo di overhead (sincronizzazione)
  rispetto al dinamico puro.
\end{itemize}

\pandocbounded{\includegraphics[keepaspectratio]{Implementazione-Parallela_files/figure-latex/unnamed-chunk-12-1.pdf}}

\subsubsection{Analisi in base all'inizializzazione delle Griglie (1D vs
2D)}\label{analisi-in-base-allinizializzazione-delle-griglie-1d-vs-2d}

L'obiettivo di questo studio è analizzare come le performance del
programma variano in base all'architettura della griglia. In particolare
si analizzano due configurazioni diverse di griglia:

\begin{itemize}
\tightlist
\item
  \textbf{Vettore monodimensionale}, dov'e la griglia'essa viene
  vettorizzata. In questo caso per spostarsi all'interno del vettore si
  usa l'indice ``y*N + x''
\item
  \textbf{Matrice bidimensionale}, dove si usa una sintassi più naturale
  dove per muoversi usiamo separatamente gli indici ``x ed y''
\end{itemize}

\pandocbounded{\includegraphics[keepaspectratio]{Implementazione-Parallela_files/figure-latex/unnamed-chunk-14-1.pdf}}

\pandocbounded{\includegraphics[keepaspectratio]{Implementazione-Parallela_files/figure-latex/unnamed-chunk-16-1.pdf}}

Questo risultato è estremamente soddisfacente, perché si può notare come
la scelta di un vettore monodimensionale, seppur pensato per ottimizzare
il comportamento della CPU riguardo i campionamenti dalla RAM, non
sempre risulti la decisione più corretta. Si può notare infatti che il
comportamento della matrice si adatta meglio, anche se di poco, al
multithreading. D'altro canto, il vettore monodimensionale risulta
sempre più performante all'aumentare della dimensione del problema.

\end{document}
